\subsection{Status evaluasi dan ruang lingkup hasil}
\label{sec:exploratory-status}

Hasil berikut berasal dari eksperimen pengembangan yang telah selesai,
bukan evaluasi konfirmatori pada himpunan uji final. Label MAL mengikuti
sumber akuisisi; kandidat HN belum memiliki verifikasi benign oleh anotator
independen. Karena itu, kami melaporkan jumlah keputusan \emph{malicious},
validitas keluaran, dan perubahan keputusan, tanpa menafsirkan alarm HN
sebagai FPR atau menghitung akurasi terhadap gold. Card KB malicious/benign
terverifikasi dan pemisahan final KB--test belum selesai. Perbandingan utama
DIFF+I+MB dengan DIFF+I+MM serta ablasi intent belum memiliki hasil
konfirmatori.

\subsection{Diagnosis bersyarat pada versi dependency}
\label{sec:dependency-diagnosis}

Eksperimen D88 meninjau 20 pembaruan yang menambahkan
\texttt{easy-day-js}. Kelompok ini dipilih setelah analisis kesalahan pada
dev set; seluruh target berbagi satu keluarga dependency dan tidak dianggap
sebagai 20 insiden independen. Qwen3-4B-Instruct-2507 menerima dua snapshot
dependency secara terpisah, menghasilkan 40 inferensi baru. Instruksi
kondisional pada kedua snapshot identik. Kontrol IDS-I tanpa konteks
diambil dari 20 respons historis D87; terhadap kontrol tersebut, penambahan
instruksi kondisional dan isi konteks tidak terpisahkan.

\begin{table}[t]
\caption{Keputusan Qwen3-4B pada 20 target D88. Angka merupakan jumlah
verdict, bukan estimasi recall. Dua baris snapshot adalah inferensi baru.}
\label{tab:dependency-versions}
\centering
\begin{tabular}{lrrr}
\toprule
Kondisi & Malicious & Benign & Invalid \\
\midrule
IDS-I historis & 0 & 20 & 0 \\
Snapshot 1.11.21 & 0 & 20 & 0 \\
Snapshot 1.11.22 & 20 & 0 & 0 \\
\bottomrule
\end{tabular}
\end{table}

Seluruh target berubah dari benign menjadi malicious ketika snapshot
1.11.21 diganti dengan 1.11.22 (Tabel~\ref{tab:dependency-versions}).
Versi 1.11.22 memiliki postinstall yang menjalankan \texttt{setup.cjs};
pembacaan statis file tersebut menunjukkan pengambilan konten jaringan,
penulisan file, pemanggilan proses anak, dan penonaktifan verifikasi TLS.
Sebanyak 18 dari 20 respons positif mengutip ID konteks dependency,
tetapi hanya 14 mengutip isi \texttt{setup.cjs}. ID kutipan yang sah belum
membuktikan bahwa seluruh alasan model didukung secara semantik.

Kedua versi memenuhi rentang \texttt{\textasciicircum1.11.21} dan
mendahului penerbitan target menurut metadata arsip. Namun, versi yang
benar-benar terpasang tidak diketahui karena tidak tersedia lockfile.
Hasil 20/20 malicious karenanya merupakan keputusan bersyarat pada snapshot,
bukan recall 100\%. Tidak ada keluaran terpotong atau ID bukti tak dikenal.
Total penggunaan 40 inferensi adalah 370.164 token prompt dan output.

\subsection{Kontrol konteks pada dependency lain}
\label{sec:dependency-generalization}

D89 memakai 15 target dengan 15 nama dependency berbeda dan mengecualikan
seluruh target \texttt{easy-day-js}. Kohort terdiri dari 14 kandidat HN dan
satu MAL menurut sumber, yaitu \texttt{zuper-cli@1.0.1}. Dari 67 kombinasi
target--versi lokal yang memenuhi semver, 44 dikeluarkan karena dependency
diterbitkan setelah target. Dua kombinasi self-dependency, satu versi lokal
lebih lama, dan lima kombinasi yang melampaui batas nama/target juga
dikeluarkan. Versi dipilih dari arsip lokal yang tersedia dan mendahului
target; prosedur ini tidak menetapkan versi instalasi aktual atau versi
terbaru dari seluruh registry. Enam target pernah terpapar pada D86;
seluruh kohort tetap merupakan data pengembangan.

Kondisi NO-SOURCE dan WITH-SOURCE memakai instruksi kondisional serta
nama/versi dependency yang sama. WITH-SOURCE menambahkan isi snapshot;
NO-SOURCE menyatakan bahwa isi tidak tersedia. Qwen3-4B menghasilkan
30 respons valid: seluruh 14 kandidat HN dan satu MAL dinilai benign pada
kedua kondisi. Tidak ada perubahan keputusan. Hanya dua respons WITH-SOURCE
mengutip isi dependency. Total token meningkat dari 67.749 pada NO-SOURCE
menjadi 128.286 pada WITH-SOURCE (1,89 kali), tanpa perubahan verdict.

Semua kandidat HN terpilih hanya memiliki sinyal perubahan dependency;
tidak ada snapshot terpilih yang memiliki install hook. Dengan demikian,
kohort ini belum mewakili pembaruan sah dengan penambahan perilaku runtime
sensitif. Nol alarm tidak menunjukkan FPR 0\%, dan satu MAL yang tetap
benign tidak mendukung klaim kenaikan deteksi. Lima belas nama dependency
juga tidak membuktikan 15 keluarga serangan independen.

\subsection{Pembanding keluarga model: Granite 3.3 8B}
\label{sec:granite-comparison}

Eksperimen D90 menjalankan IBM Granite 3.3 8B Instruct pada seluruh 15
target dan dua kondisi D89, menghasilkan 30 inferensi baru. Teks masukan,
snapshot, schema JSON, parser, dan urutan permintaan dipertahankan.
Pembanding Qwen menggunakan 30 respons historis D89. Kedua model memakai
BF16, temperature 0, top-p 1, seed 0, batas output 2.048 token, konteks
32.768 token, dan satu permintaan pada satu waktu. Granite menggunakan
template native dengan \texttt{thinking=false} dan token akhir native;
prompt maksimum 24.467 token, tanpa pemotongan tambahan. Bobot model dan
input dipin menggunakan SHA-256 sebelum inferensi. Perbedaan keluarga,
ukuran, tokenizer, dan template berarti perbandingan ini tidak mengisolasi
efek jumlah parameter.

\begin{table*}[t]
\caption{Perbandingan pada target dan kondisi identik: Qwen3-4B (historis
D89) dan Granite 3.3 8B (baru D90). Sel menunjukkan malicious/benign/invalid.
HN belum terverifikasi benign; MAL mengikuti label sumber.}
\label{tab:granite-comparison}
\centering
\begin{tabular}{llrrr}
\toprule
Kohort & Kondisi & Target & Qwen3-4B & Granite 3.3 8B \\
\midrule
HN pending & NO-SOURCE & 14 & 0/14/0 & 0/14/0 \\
HN pending & WITH-SOURCE & 14 & 0/14/0 & 0/14/0 \\
MAL sumber & NO-SOURCE & 1 & 0/1/0 & 0/1/0 \\
MAL sumber & WITH-SOURCE & 1 & 0/1/0 & 0/1/0 \\
\bottomrule
\end{tabular}
\end{table*}

Tidak ada satu pun verdict yang berbeda antar model
(Tabel~\ref{tab:granite-comparison}). Seluruh respons Granite memenuhi
schema, tanpa keluaran terpotong, ID bukti tak dikenal, atau error transport.
Satu target MAL tetap benign, sehingga pergantian model belum memperbaiki
deteksi pada kohort ini. Dua respons WITH-SOURCE mengutip isi dependency,
pada \texttt{get-symbol-description@1.0.2} dan
\texttt{@tolicodes/picklejs@0.75.0}; dukungan semantik alasan belum
diverifikasi manusia.

Granite memakai 76.964 token pada NO-SOURCE dan 146.418 token pada
WITH-SOURCE (total 223.382), dibandingkan total 196.035 token Qwen.
Token antar model mengikuti tokenizer masing-masing dan tidak disamakan
dengan biaya komputasi. Median latensi pada 14 kandidat HN adalah 1.676 ms
dan 2.422 ms untuk dua kondisi Granite; angka historis Qwen adalah
872 ms dan 1.195 ms. Latensi merupakan observasi sesi dengan prefix caching,
tanpa warmup generasi, bukan benchmark kecepatan terkontrol. Audit
pasca-inferensi mencocokkan payload, jumlah token, parsing, kutipan,
27 checksum sumber, enam artefak kontrol, dan 14 berkas model. Hasil ini
mendukung kesamaan keputusan pada kohort yang diuji, tanpa memberikan
peringkat umum model atau bukti penurunan FPR.
